% G07 - Fase 2, Entrega 03 - apresentacao
\documentclass[aspectratio=169,11pt]{beamer}
\usetheme{default}
\usepackage{fontspec}
\setsansfont{TeX Gyre Heros}
\usepackage{booktabs,amsmath,tikz,array}
\usepackage{newunicodechar}
\newunicodechar{Ω}{\ensuremath{\Omega}}
\newunicodechar{→}{\ensuremath{\to}}
\newunicodechar{×}{\ensuremath{\times}}
\newunicodechar{·}{\ensuremath{\cdot}}
\newunicodechar{≈}{\ensuremath{\approx}}
\newunicodechar{≤}{\ensuremath{\le}}
\newunicodechar{°}{\ensuremath{^\circ}}
\newunicodechar{µ}{\ensuremath{\mu}}
\newunicodechar{σ}{\ensuremath{\sigma}}
\newunicodechar{−}{\ensuremath{-}}
\definecolor{az}{HTML}{1F4E8C}\definecolor{azc}{HTML}{E8F0FA}
\definecolor{vm}{HTML}{C0392B}\definecolor{cz}{HTML}{5A6270}
\setbeamercolor{normal text}{fg=black!85}
\setbeamercolor{structure}{fg=az}
\setbeamercolor{frametitle}{fg=az}
\setbeamerfont{frametitle}{size=\Large,series=\bfseries}
\setbeamertemplate{navigation symbols}{}
\setbeamertemplate{itemize item}{\color{az}\rule[0.25ex]{0.6ex}{0.6ex}}
\setbeamertemplate{frametitle}{\vspace{2pt}\insertframetitle\par\vspace{1pt}{\color{az}\rule{2.2cm}{2pt}}\vspace{-4pt}}
\setbeamertemplate{footline}{\leavevmode\hbox to\paperwidth{\hspace{12pt}\scriptsize\color{cz}\rodape\hfill\insertframenumber\hspace{12pt}}\vspace{6pt}}
\newcommand{\fig}[2][\textwidth]{\includegraphics[width=#1]{../relatorio/figuras/#2}}
\newcommand{\kpi}[3][2.55cm]{\begin{tikzpicture}\node[fill=azc,rounded corners=3pt,inner sep=6pt,text width=#1,align=center]{{\Large\bfseries\color{az}#2}\\[1pt]{\scriptsize\color{cz}#3}};\end{tikzpicture}}
\newcommand{\kpivm}[3][2.55cm]{\begin{tikzpicture}\node[fill=vm!10,rounded corners=3pt,inner sep=6pt,text width=#1,align=center]{{\Large\bfseries\color{vm}#2}\\[1pt]{\scriptsize\color{cz}#3}};\end{tikzpicture}}
\newcommand{\capa}[3]{{\setbeamertemplate{footline}{}\setbeamercolor{background canvas}{bg=az}
\begin{frame}\color{white}\vspace{1.4cm}
{\small #1}\\[10pt]{\huge\bfseries #2}\\[6pt]{\Large #3}\\[24pt]
{\color{white}\rule{3cm}{1.5pt}}\\[10pt]
\small Nícolas De Almeida Souza · Lucas Momesso Alves · Cláudio Henrique da Costa Santos
\end{frame}}}
\newcommand{\destaque}[1]{\begin{center}\colorbox{azc}{\parbox{.9\textwidth}{\centering\vspace{5pt}#1\vspace{5pt}}}\end{center}}
\newcommand{\pontos}[6]{\begin{tabular}{@{}>{\raggedright\arraybackslash}p{.3\textwidth}@{\hspace{.035\textwidth}}>{\raggedright\arraybackslash}p{.3\textwidth}@{\hspace{.035\textwidth}}>{\raggedright\arraybackslash}p{.3\textwidth}@{}}
#1 & #3 & #5\\[4pt]\small #2 & \small #4 & \small #6\end{tabular}}
\newcommand{\contexto}[1]{{\small\color{cz}#1}\par\vspace{4pt}}
\newcommand{\rodape}{Instrumentação em malha aberta}
\begin{document}

\capa{Fase 2 · Entrega 03}{Instrumentação em malha aberta}{Medindo a corrente de fase do motor BLDC}

\begin{frame}{O que o instrumento precisa medir}
\contexto{A corrente de fase simulada no VHIL (forma de 120°, 244 Hz, ondulação de PWM de 10 kHz) define os requisitos do instrumento.}
\fig[.9\textwidth]{fig_demanda.pdf}\\[4pt]
\kpi[3.05cm]{±10 A}{faixa (partida 6,4 A)}\hfill
\kpi[3.05cm]{3,5 kHz}{banda (13.ª harmônica)}\hfill
\kpi[3.05cm]{≤ 25 mA}{viés (2 \% do nominal)}\hfill
\kpi[3.05cm]{≤ 100 µs}{atraso (1 ciclo PWM)}
\end{frame}

\begin{frame}{Sensor escolhido: efeito Hall ACS712-20A}
\contexto{Entre o shunt e o Hall, o Hall é o mais simples de ligar ao inversor: isolado e já com saída em tensão.}
\begin{columns}[c]
\column{.52\textwidth}
\begin{itemize}\setlength{\itemsep}{12pt}
\item \textbf{Isolado} do circuito de potência
\item \textbf{Saída em tensão}, sem amplificador externo
\item Versão de \textbf{20 A}: a de 5 A saturaria na partida
\item Base: Aguirre, seções 7.6 e 8.2.1
\end{itemize}
\column{.4\textwidth}\centering
\kpi{100 mV/A}{sensibilidade}\\[6pt]
\kpi{80 kHz}{banda}\\[6pt]
\kpi{±1,5 \%}{erro total}
\end{columns}
\end{frame}

\begin{frame}{A cadeia de medição no Typhoon HIL}
\contexto{Cada bloco é uma C function com uma característica do sensor ou do condicionamento, ajustável pelo SCADA. O caminho ideal é um ganho 1.}
\fig[.82\textwidth]{fig_cadeia.pdf}\\[0pt]
\begin{columns}[T]
\column{.46\textwidth}
{\color{vm}\bfseries Imperfeições do sensor}
\begin{itemize}\footnotesize
\item Erros de ganho e de zero
\item Temperatura de 60 °C
\item Banda finita e ruído
\end{itemize}
\column{.46\textwidth}
{\color{az}\bfseries Condicionamento}
\begin{itemize}\footnotesize
\item Filtro RC de 15 kHz
\item ADC de 12 bits síncrono com o PWM
\item Correção calculada + média por período
\end{itemize}
\end{columns}
\end{frame}

\begin{frame}{A medição real atende à demanda}
\contexto{Três ensaios no VHIL (partida, corrente nula e pulso de carga) comparam a medição ideal com a real.}
\begin{columns}[c]
\column{.66\textwidth}\fig{fig_ensaios.pdf}
\column{.28\textwidth}\centering
\kpi{≤ 2 mA}{viés}\\[6pt]
\kpi{58 mA}{erro rms (5 \%)}\\[6pt]
\kpi{60 µs}{atraso}
\end{columns}
\end{frame}

\begin{frame}{O desvio de zero domina o viés}
\contexto{Para saber qual característica pesa mais, cada bloco também calcula a medição só com a sua imperfeição.}
\fig[.8\textwidth]{fig_contribuicoes.pdf}\\[0pt]
\begin{columns}[c]
\column{.5\textwidth}\footnotesize
O zero desloca o valor médio e não sai com média.\\[4pt]
Ruído e ADC pesam no erro instantâneo, que a média por período reduz.
\column{.46\textwidth}\centering
\kpivm{75,5 mA}{sem correção, com $i = 0$}\hspace{4pt}\kpi{1,2 mA}{com correção}
\end{columns}
\end{frame}

\begin{frame}{Exagerando as características}
\contexto{Banda, ruído e bits do ADC foram levados a níveis extremos, com o erro previsto por fórmula antes de cada simulação.}
\begin{columns}[c]
\column{.64\textwidth}\fig{fig_exacerbacao.pdf}
\column{.29\textwidth}
\begin{itemize}\setlength{\itemsep}{10pt}\footnotesize
\item Previsões feitas \textbf{antes} de simular batem com o resultado
\item Exceção: ADC de 4 bits, em que o degrau (3,1 A) é maior que o sinal
\item Ruído de 1 V gera \textbf{10 A rms} de erro
\end{itemize}
\end{columns}
\end{frame}

\begin{frame}{Conclusões}
\contexto{O que cada característica significa para o motor BLDC em operação.}
\vspace{14pt}
\pontos{{\color{vm}\Large\bfseries Zero}}{O mais nocivo em uso real: viés sistemático no torque e na potência.}{{\color{vm}\Large\bfseries Ruído}}{O pior quando exagerado: disparos falsos da proteção.}{{\color{az}\Large\bfseries Banda}}{Banda baixa atrasa a medida e apaga a comutação.}
\vspace{22pt}
\destaque{Com correção calculada e média por período, o ACS712 atende a toda a demanda.}
\end{frame}

\end{document}
